\section{Keep simulated work separate from encoded size}
The first implementation maps signed differences with shifts and XOR and counts output
groups with a loop. The alternative branches on sign and uses size thresholds. Its
decoder accumulates weighted base-128 digits with a multiplier and branches for inverse
mapping. Both satisfy the same byte-level contract and receive the same correctness
checks.

\begin{table}[htbp]
\centering
\caption{Whole-function retired instructions for the checked encoder prototypes.}
\tagpdfsetup{table/header-rows={1}}
\begin{tabular}{lrr}
\toprule
Dataset & Shift-count encoder & Threshold encoder\\
\midrule
Steady & 1041 & 1072\\
Ramp & 1026 & 1062\\
Timestamps & 1551 & 1546\\
Spike & 1086 & 1120\\
Alternating endpoints & 2946 & 2850\\
Random & 2931 & 2840\\
\bottomrule
\end{tabular}
\end{table}

The interval begins at function entry and stops before its supplied continuation.
It includes frame setup, capacity checks, all complete values and restoration. It
excludes assembly, session creation, caller storage and I/O. The decoder measurements
use independently modeled encoded bytes, rather than charging the encoder inside the
decoder's interval. The two implementations have identical output sizes while their
instruction totals differ by dataset. These counts are not elapsed host time or cycles.

Possible experiments include avoiding a second scan of the mapped value, retaining a
short candidate in a private buffer, or checking a capacity upper bound before exact
length calculation. Every strategy still needs item-atomic output and the specified
zero-capacity behavior. Include private preparation work in its complete-call total.
A five-byte worst-case reservation cannot replace the exact capacity check: a one-byte
value must succeed when only one byte remains.
